\documentclass[12pt]{article}
\usepackage[a4paper,margin=1in]{geometry}
\usepackage{amsmath}
\usepackage{amssymb}
\usepackage{enumitem}
\usepackage[T1]{fontenc}

\begin{document}
\pagenumbering{gobble}

\begin{center}
{\LARGE\bfseries Answer Key}\\[0.3em]
{\large Challenge: Decimals, Percentages, and Fractions}
\end{center}

\begin{enumerate}[label=\textbf{Q\arabic*.}]

\item
\begin{enumerate}[label=(\alph*)]
\item $0.125=\dfrac{1}{8}=12.5\%$.
\item $0.203125=\dfrac{13}{64}=20.3125\%$.
\end{enumerate}

\item
\begin{enumerate}[label=(\alph*)]
\item $\dfrac{7}{20}=0.35=35\%$.
\item $\dfrac{33}{320}=0.103125=10.3125\%$.
\end{enumerate}

\item
\begin{enumerate}[label=(\alph*)]
\item $17.5\%=0.175=\dfrac{7}{40}$.
\item $57.8125\%=0.578125=\dfrac{37}{64}$.
\end{enumerate}

\item
\textbf{Basis points.}
\begin{enumerate}[label=(\alph*)]
\item $1\text{ bp}=0.01\%=0.0001=\dfrac{1}{10000}$.
\item $175\text{ bp}=1.75\%=0.0175=\dfrac{7}{400}$.
\end{enumerate}

\medskip
\textbf{Gilts.}
\begin{enumerate}[label=(\alph*), start=3]
\item $5.678\%=0.05678=\dfrac{2839}{50000}$.
\item Annual coupon: \pounds 56.68? \\[The source says 5.678\% of 1000 = 56.78, so use that:]

Wait: annual coupon is $56.78$; coupons total $30(56.78)=1703.40$; including repayment, total cost is \pounds 2703.40.

\medskip
\item $43.5\text{ bp}=0.435\%=0.00435$; the coupon rises by \pounds 4.35 per year, so the extra cost is $30(4.35)=\pounds 130.50$.

\item Compounding can be achieved by reinvesting every coupon payment and maturing principal in further gilts as soon as they are received, so interest is earned on previous interest.
\end{enumerate}

\newpage

\item
\begin{enumerate}[label=(\alph*)]
\item Upward-sloping, since \[
4.377\%<4.474\%<4.488\%<4.492\%<4.544\%.
\]
\item The 5-year gilt has the highest yield, so it gives the best simple return if held to maturity and is the best starting point for a reinvestment strategy.
\item Use a coupon-reinvestment ladder: start with the 5-year gilt and each year reinvest all coupon cash received at the previous year-end in the longest-dated gilt that still matures by the end of Year 4.
\item Using that ladder, the final value is \pounds 12,484.60.
\[
\begin{array}{c|c}
\text{Gilt} & \text{Maturity payment}\\ \hline
5\text{-year} & 10454.40\\
4\text{-year} & 474.81\\
3\text{-year} & 496.11\\
2\text{-year} & 518.30\\
1\text{-year} & 540.98
\end{array}
\]
\[
10454.40+474.81+496.11+518.30+540.98=12484.60.
\]
\item Annual simple rate: $r_s=0.049692=4.9692\%$.
\item Annual compound rate: $r_c\approx 0.045382=4.538\%$.
\item Monthly rate: $r_m\approx 0.003705=0.3705\%$.
\end{enumerate}

\end{enumerate}

\end{document}